% ============================================================
%  Глава 3 — Откуда берётся состояние
% ============================================================
\chapter{Откуда берётся состояние}
\label{ch:state}

\section{Вопрос}

В \eqref{eq:world_kinematics} мы писали $\dot{p}_x = v\cos\theta$. Это
дифференциальное уравнение --- то, как состояние \emph{меняется во времени}.
Но MPC, который мы будем запускать $50$~раз в секунду, ему на каждом тике
нужно $\vect{x}_k$ --- \emph{измеренное} состояние \emph{сейчас}.

Откуда оно? Робот же не «знает» своих координат --- их нужно посчитать
из сигналов реальных датчиков.

\begin{ideabox}
Это разделение --- очень важное:
\begin{itemize}
  \item \emph{Модель} говорит, как состояние должно меняться (формула).
  \item \emph{Оценка состояния} (state estimation) даёт текущее значение
        состояния по показаниям датчиков (число).
\end{itemize}
В этой главе нас интересует \emph{оценка}. Модель будет дальше.
\end{ideabox}

\section{Состояние в МПС-сценарии}

В сценарии «D метров вперёд» состояние --- это пять чисел:

\begin{formulabox}[title={Канонический МПС-вектор состояния}]
\begin{equation}
  \vect{x} \;=\; \begin{pmatrix} s \\ v \\ \theta \\ \omega \\ e_\text{int} \end{pmatrix}
  \;\in\; \Real^5
  \label{eq:mps_state}
\end{equation}

\begin{tabular}{lll}
\toprule
$s$           & пройденная дистанция         & м\\
$v$           & продольная скорость           & м/с\\
$\theta$      & курс (yaw)                    & рад\\
$\omega$      & угловая скорость              & рад/с\\
$e_\text{int}$& интеграл ошибки скорости      & м ($\approx v\cdot t$)\\
\bottomrule
\end{tabular}
\end{formulabox}

\begin{notebox}
\textbf{Почему именно так, а не $(p_x, p_y, \theta, v, \omega)$?}

В общей литературе по дифференциальным роботам состояние ---
$(p_x, p_y, \theta, v, \omega)$ (две декартовы координаты + курс +
две скорости). И в проекте есть «легаси»-форма с такими пятью переменными
(\filepath{pi\_nodes/control/state\_space\_model.py}: переменные
\texttt{px, py, theta, v, omega}). Но для сценария «D метров вперёд»
эта форма \emph{избыточна}: при правильной работе $p_y \equiv 0$ (робот
едет прямо), а $p_x$ совпадает с $s$.

Поэтому в МПС-сценарии мы переходим к более компактной форме
\eqref{eq:mps_state}: одна продольная координата $s$ вместо двух,
плюс \emph{интегральный член} $e_\text{int}$, который ловит установившуюся
ошибку по скорости (без него робот хронически не добирает $v_\text{target}$).
\end{notebox}

\begin{warnbox}
В коде live ОБЕ модели существуют:
\begin{itemize}
  \item Легаси $(p_x, p_y, \theta, v, \omega)$ --- секция \texttt{control.}
        в \filepath{config.yaml}, используется старым PID-контуром.
  \item Каноническая $(s, v, \theta, \omega, e_\text{int})$ --- секция
        \texttt{mps:}, используется в MPS-сценарии.
\end{itemize}
Это \emph{не баг} --- это разделение «учебной» подсистемы и боевой.
Просто когда читаешь код, всегда смотри, \emph{какая} секция
(\texttt{control.*} или \texttt{mps.*}) грузит матрицы.

Историческая ловушка: \texttt{mps\_node} в коммите \texttt{8521742}
случайно строился из \texttt{control.*} и трактовал свой
$x[1] = v$ как $p_y$, после чего MPC доворачивал робота, пытаясь
свести $p_y$ к $v_\text{target}$. Робот ездил по кругу.
\end{warnbox}

\section{Источники измерений}

Каждая компонента \eqref{eq:mps_state} приходит не «магически», а из
конкретного сенсора через конкретный фильтр.

\begin{center}
\begin{tabular}{llp{6cm}}
\toprule
\textbf{Компонента} & \textbf{Сенсор} & \textbf{Преобразование}\\
\midrule
$s, v$    & энкодеры моторов   & интегрирование числа тиков $\to$ путь;
                                 \texttt{velocity\_ekf}\\
$\theta, \omega$ & MPU-6050 IMU (гироскоп) & \texttt{ekf\_imu}: фильтр Калмана
                                            1D на $\theta$\\
$e_\text{int}$ & --- & накапливается \emph{самим} \texttt{mps\_node} как
                       $\sum (v - v_\text{target})\,\Delta t$\\
\bottomrule
\end{tabular}
\end{center}

Подробное описание фильтров --- в общем справочнике
\filepath{latex\_doc/chapters/01\_velocity\_ekf.tex},
\filepath{latex\_doc/chapters/02\_imu\_ekf.tex}.

\subsection{Одометрия по гусеницам}

Энкодеры считают, сколько раз мотор провернулся за последний тик.
Зная диаметр приводного колеса и передаточное число, это превращается
в пройденный гусеницей путь $\Delta s_L, \Delta s_R$. Скорости и
позиция центра:

\begin{equation}
  v = \frac{\Delta s_L + \Delta s_R}{2\,\Delta t},
  \qquad
  \omega_\text{odom} = \frac{\Delta s_L - \Delta s_R}{B\,\Delta t},
  \qquad
  s_{k+1} = s_k + v\,\Delta t.
\end{equation}

Эти формулы --- та же \eqref{eq:diff_inverse}, что в главе~\ref{ch:kinematics},
только применённая к \emph{измеренным} $\Delta s$, а не к команде.

\subsection{Курс по IMU}

Гироскоп MPU-6050 даёт измеренную угловую скорость $\omega_\text{imu}(t)$.
Интегрируя её, получим курс $\theta(t) = \int \omega_\text{imu}\,dt$.
Проблема: гироскоп дрейфует. За минуту накопится несколько градусов
ошибки даже на стоящем роботе.

Для коротких сценариев ($D \le 5$~м $\Rightarrow$ $t \le 25$~с при
$v_\text{target} = 0.2$~м/с) дрейф пренебрежимо мал, и мы используем
прямое интегрирование с лёгкой EKF-сглаживанием. Для длинных сессий
проект делает дополнительную привязку через \texttt{position\_fusion}
(акселерометр + одометрия), но в МПС-сценарии это не критично.

\section{Слияние: \texttt{position\_fusion}}

Все источники сводятся в одну ноду \texttt{position\_fusion}, которая
публикует MQTT-сообщение \texttt{odom} с полями \texttt{x, vx, theta, vz}.
Это и есть «питающие» измерения для \texttt{mps\_node}.

\begin{codebox}[title={\filepath{pi\_nodes/nodes/mps\_node.py:342}}]
\begin{lstlisting}[language=Python]
def _on_odom(self, topic: str, payload):
    # samurai/{id}/odom — motor_node публикует:
    #   x в САНТИМЕТРАХ, vx в м/с, theta в рад, vz в рад/с.
    s = float(payload.get('x', 0.0)) / 100.0   # см -> м
    v = float(payload.get('vx', 0.0))
    theta = float(payload.get('theta', 0.0))
    omega = float(payload.get('vz', 0.0))
    with self._lock:
        self._x_meas = np.array(
            [s, v, theta, omega, self._x_meas[_EINT]]
        )
        self._x_meas_ts = time.time()
\end{lstlisting}
\end{codebox}

\begin{warnbox}
\textbf{Подводный камень --- единицы.} \texttt{motor\_node} публикует
позицию в \emph{сантиметрах}, а МПС-модель работает в \emph{метрах}.
В коммите \texttt{5f54c54} баг был такой: эвристика «если $|s| > 20$,
то это уже метры, не конвертируй» оставляла первые 20~см не
сконвертированными --- и MPC видел фантомную ошибку позиции $\times 100$,
прижимая $u$ к насыщению $u_\text{max}$. Исправлено безусловной
конвертацией.

Урок: в любой межкомпонентной границе \emph{явно фиксируй единицы}.
\end{warnbox}

\section{Старт сценария: «обнуление» начала}

Когда пользователь жмёт «Run», \texttt{mps\_node} запоминает \emph{текущие}
$s$ и $\theta$ как $s_\text{start}, \theta_\text{start}$ (см. \texttt{\_RunState}).
Дальше состояние для MPC считается \textbf{относительно старта}:

\begin{equation}
  x[\,s\,]_\text{rel}     = s - s_\text{start},
  \qquad
  x[\,\theta\,]_\text{rel} = \text{wrap}_{[-\pi, \pi]}(\theta - \theta_\text{start}).
  \label{eq:relative_state}
\end{equation}

\begin{ideabox}
Зачем? Сценарий говорит: «проехать $D$ метров \emph{отсюда}». Если бы
мы кормили в MPC \emph{абсолютную} $s$, то reference $s_\text{ref}$
тоже надо было бы делать абсолютным --- а робот мог стартовать с
$s=4.7$~м после предыдущего прогона. Гораздо проще обнулить отсчёт
в момент старта.

То же с курсом: $\theta_\text{ref} = \varphi$ означает «отвернуть на
$\varphi$ от \emph{текущего} направления», а не от мирового севера.
\end{ideabox}

Функция \texttt{\_normalize\_angle} оборачивает разность в $[-\pi, \pi]$,
чтобы поворот на $-179^\circ$ не превратился в $+181^\circ$ (что бы
сломало MPC, который квадратично штрафует ошибку курса).

\section{Что мы получили}

\begin{itemize}
  \item Состояние МПС-сценария --- пять чисел $(s, v, \theta, \omega, e_\text{int})$.
  \item Первые четыре приходят из MQTT-топика \texttt{odom} с Pi-ноды
        \texttt{position\_fusion}, конвертируясь см$\to$м.
  \item Пятый ($e_\text{int}$) живёт внутри самого \texttt{mps\_node}
        --- это интеграл ошибки скорости, нужен для устойчивого
        отслеживания $v_\text{target}$.
  \item На старте прогона $s$ и $\theta$ \emph{обнуляются}: дальнейшие
        вычисления идут в относительных координатах.
\end{itemize}

Теперь у нас есть и состояние, и формула \eqref{eq:world_kinematics}, как
оно меняется кинематически. Но кинематики мало: моторы инерционны,
команда отрабатывается не мгновенно. Это --- следующая глава.
